\section{语言模型评估：困惑度}

\subsection{什么是困惑度}

语言模型的本质功能是为一段文本分配概率。评估语言模型的标准方法是：给定一段人类书写的\textbf{未见过的文本}（非训练数据），让模型逐词预测，看模型预测得有多好。

\begin{importantbox}{困惑度的定义}
困惑度（Perplexity）是交叉熵的指数形式。设语言模型对测试文本 $w_1, w_2, \ldots, w_T$ 的预测概率为 $P(w_t | w_1, \ldots, w_{t-1})$，则：

$$\text{Perplexity} = \prod_{t=1}^{T} \left( \frac{1}{P(w_t | w_1, \ldots, w_{t-1})} \right)^{1/T}$$

等价地，困惑度 = $\exp(\text{Cross Entropy})$，即：
$$\text{Perplexity} = \exp\left( -\frac{1}{T} \sum_{t=1}^{T} \log P(w_t | w_1, \ldots, w_{t-1}) \right)$$
\end{importantbox}

\begin{figure}[H]
\centering
\includegraphics[width=0.95\textwidth]{slides-images/slide-002.jpg}
\caption{语言模型评估：Perplexity 的定义与直觉解释\protect\footnotemark}
\end{figure}
\footnotetext{Slides 第3页。}

\subsection{困惑度的直觉理解}

困惑度可以理解为\textbf{等效均匀选择数}：如果困惑度为 64，就相当于每次从 64 个等概率的选项中猜测下一个词。困惑度越低，模型越好。

\begin{knowledgebox}{困惑度的历史由来}
困惑度这一概念由 IBM 的 Fred Jelinek 在 1970--80 年代引入。当时，符号主义 AI 的研究者对信息论和交叉熵不甚了解，Jelinek 需要一个更直观的度量来解释统计语言模型的效果。困惑度通过指数变换将交叉熵转化为``等效选择数''，更容易被非数学背景的人理解。尽管从现代视角看，直接使用交叉熵（负对数似然）更为自然，但困惑度作为评估指标已经约定俗成。
\end{knowledgebox}

\begin{warningbox}{比较困惑度时注意对数底数}
困惑度的值取决于对数底数的选择。传统上使用以 2 为底的对数（信息论传统），现代实践中更多使用自然对数。比较不同论文的困惑度数值时，务必确认所用的对数底数一致，否则数值不可直接比较。
\end{warningbox}

\subsection{语言模型的历史进展}

\begin{figure}[H]
\centering
\includegraphics[width=0.95\textwidth]{slides-images/slide-003.jpg}
\caption{语言模型困惑度的历史进展：从 n-gram 到 LSTM 再到现代大语言模型\protect\footnotemark}
\end{figure}
\footnotetext{Slides 第4页。}

\begin{table}[H]
\centering
\begin{tabular}{lcc}
\toprule
\textbf{模型} & \textbf{困惑度} & \textbf{时期} \\
\midrule
Interpolated Kneser-Ney (5-gram) & $\sim$67 & 2000s \\
RNN + MaxEnt 混合 & $\sim$51 & 2010s 初期 \\
LSTM 语言模型 & 30--43 & 2015--2018 \\
现代大语言模型 & 个位数 & 2020s \\
\bottomrule
\end{tabular}
\caption{不同时期语言模型在标准测试集上的困惑度}
\end{table}

从 n-gram 时代的困惑度 $\sim$67，到早期 RNN 的 $\sim$51，再到 LSTM 时代的 30--43，困惑度减半意味着交叉熵降低约 1 bit。现代大语言模型的困惑度已降至个位数，意味着模型大多数时候能准确预测下一个词。

\subsection{本章小结}

困惑度是语言模型的标准评估指标，本质上是交叉熵的指数形式。数值越低代表模型越好。从 n-gram 到 LSTM 再到 Transformer，语言模型的困惑度经历了数量级的改善。

%% ========== 梯度问题 ==========

